This chapter reviews the literature relevant to language-model-based task planning in heterogeneous multi-robot systems. It examines the evolution of task planning approaches, recent advances in language-model-based planning, grounding and environment representation, language model selection, and simulation environments. Together, these developments establish the current state of the field and motivate the research presented in this thesis.

Task planning is a fundamental problem in multi-robot systems, concerned with allocating, ordering, and coordinating tasks among multiple robots to achieve a shared objective efficiently. Unlike single-robot planning, multi-robot systems must additionally consider interactions between agents, including resource sharing, communication, task dependencies, and conflicting objectives. As the number of robots and the complexity of the environment increase, the planning problem becomes increasingly difficult due to the combinatorial growth of possible task allocations and execution sequences.

Early research addressed these challenges using explicitly engineered planning and coordination strategies based on optimisation, heuristics, and predefined decision rules. Gerkey and Matari\'c~\cite{gerkey2004formal} formalised the Multi-Robot Task Allocation (MRTA) problem and proposed a taxonomy for categorising task allocation methods according to robot capabilities, task characteristics, and allocation strategies. Their work established a common framework for analysing and comparing multi-robot planning algorithms and continues to influence subsequent research.

Yan \textit{et al.}~\cite{yan2013mrscoordination} subsequently reviewed the broader field of multi-robot coordination, covering auction-based task allocation, cooperative transport, communication-aware coordination, motion planning, and early learning-based methods. Representative systems such as MURDOCH demonstrated that distributed auction mechanisms could efficiently allocate tasks among heterogeneous robots, while other approaches incorporated temporal dependencies, spatial constraints, and robot availability during cooperative task execution. These studies showed that successful multi-robot task planning depends on task allocation, communication, coordination, synchronisation, and conflict resolution.

Although heuristic and optimisation-based approaches proved effective in structured environments, they relied heavily on handcrafted rules, explicit environment models, and manually designed objective functions. Adapting these systems to new environments or operational requirements often required substantial redesign, which encouraged increasing interest in learning-based approaches.

Learning-based planning, particularly Reinforcement Learning (RL) and Multi-Agent Reinforcement Learning (MARL), enables robots to learn coordination policies through repeated interaction with the environment instead of relying on predefined decision rules. Compared with classical planning methods, MARL offers greater adaptability to dynamic environments by allowing cooperative behaviours to emerge through experience. It has therefore become a prominent approach for heterogeneous multi-robot coordination.

Papoudakis \textit{et al.}~\cite{papoudakis2021benchmarking} addressed the lack of standardised evaluation in MARL by comparing representative independent learning, policy-gradient, and value-decomposition algorithms across a diverse collection of cooperative tasks. Their study showed that algorithm performance depends strongly on task characteristics such as reward sparsity and coordination complexity while introducing the EPyMARL framework and additional benchmark environments to support reproducible evaluation.

Later work by Krnjaic \textit{et al.}~\cite{krnjaic2024warehouse} applied hierarchical MARL to heterogeneous warehouse logistics involving autonomous mobile robots and human workers. Their manager-worker architecture decomposed decision making into hierarchical levels, improving sample efficiency and warehouse throughput compared with conventional MARL methods and industrial heuristic approaches. The authors also introduced the Task Assignment Robotic Warehouse (TA-RWARE) benchmark, extending RWARE to support heterogeneous task allocation and providing a reproducible platform for evaluating high-level coordination strategies.

Subsequent research shifted attention towards improving cooperation through reward design. Niu and Gürdür Broo~\cite{niu2025rewardshaping} proposed a biologically inspired reward shaping framework based on ecological symbiosis, demonstrating improved convergence stability and cooperative behaviour in heterogeneous robotic systems. Niu \textit{et al.}~\cite{niu2025symbiosis} later incorporated symbiotic principles into a Centralised Training with Decentralised Execution (CTDE) MARL framework, enabling robots to coordinate battery management and workload distribution through shared state information. Their results demonstrated improvements in task performance, resource utilisation, and energy management in simulated warehouse environments.

Despite these advances, learning-based approaches remain dependent on environment-specific training, carefully designed reward functions, and computationally intensive policy optimisation. Adapting trained policies to new environments or task configurations often requires additional training or redesign. These limitations have led researchers to investigate pretrained Large Language Models (LLMs) as an alternative paradigm for high-level task planning, seeking to exploit the reasoning capabilities acquired during large-scale pretraining instead of learning task-specific policies from scratch.

Li et al.~\cite{li2026largelanguagemodelsmultirobot} classify LLM applications in robotics according to the level of decision making they support, including high-level task planning and task allocation, mid-level motion planning, and low-level action generation. High-level planning concerns reasoning over task objectives and coordination, whereas lower levels focus on trajectory generation or direct robot control.

Existing LLM-based planning frameworks commonly employ language models as high-level decision makers while retaining conventional perception, navigation, and control modules. SayCan grounds language-model action selection with affordance estimates so that plans are biased towards actions that are both useful for the instruction and feasible in the current robotic state~\cite{ahn2022do}. Inner Monologue extends this pattern by incorporating feedback from the environment into language-model planning, while ProgPrompt represents situated task plans as program-like structures over available robot actions~\cite{huang2022innermonologue,singh2023progprompt}. Code as Policies instead uses the language model to generate executable programs over predefined robot APIs, making code the intermediate policy representation between language instructions and embodied control~\cite{liang2022code}. Related zero-shot planning work shows that language models can propose action sequences from textual task descriptions when they are given an explicit action interface, although such plans still require grounding, validation, and replanning when deployed in physical or simulated environments~\cite{huang2022lmp}.

Extending \glspl{llm} from single-robot planning to multi-robot systems introduces additional challenges related to communication, coordination, and information sharing among multiple autonomous agents. Unlike single-robot systems, where a single planning context is sufficient, multi-robot systems must determine how planning responsibilities and information should be distributed across multiple agents. Consequently, recent research has explored a range of planning and communication architectures for integrating \glspl{llm} into collaborative multi-robot systems~\cite{li2026largelanguagemodelsmultirobot}.

One line of research investigates how individual language models should be orchestrated across multiple agents. Liu \textit{et al.}~\cite{liu2023bolaa} analysed several architectures for LLM-augmented autonomous agents, ranging from simple zero-shot planning to iterative self-reflection frameworks that incorporate previous observations and actions into subsequent reasoning steps. Their work also proposed BOLAA, an orchestration strategy in which a controller manages communication among multiple specialised language-model agents. Although their study primarily considers general autonomous-agent benchmarks rather than physical robot teams, the proposed orchestration strategies provide useful insights for multi-robot coordination.

In robotics, Mandi \textit{et al.}~\cite{mandi2024roco} proposed RoCo, a multi-robot collaboration framework in which language-model agents discuss task strategies, generate sub-task plans, and revise plans using feedback such as collision checking. This work demonstrates how dialogue and shared feedback can be used to coordinate multiple robots, but it also illustrates the need for explicit grounding and validation when language models are used inside embodied planning loops.

Although \glspl{llm} possess extensive world knowledge acquired during pretraining, they do not have direct knowledge of the robot's current environment. Consequently, generated plans may reference unavailable objects, infeasible actions, or incorrect assumptions about the world state, a phenomenon commonly referred to as hallucination. Grounding addresses this problem by constraining language-model reasoning using information obtained from the environment, ensuring that generated plans remain both logically and operationally feasible.

Existing work has explored several approaches for grounding LLMs in robotic task planning. SayCan constrains language-model reasoning through affordance-based action selection, ensuring that generated actions are executable in the current environment~\cite{ahn2022do}. Feedback-based approaches further improve planning reliability by incorporating observations from the environment during execution, allowing the language model to identify infeasible plans and refine subsequent decisions through iterative replanning~\cite{huang2022innermonologue,bhat2025closedloopgrounding}. Recent surveys identify grounding as a fundamental requirement for improving the robustness and reliability of LLM-based robotic planning systems~\cite{li2026largelanguagemodelsmultirobot}.

Regardless of the grounding mechanism employed, the language model ultimately reasons over an abstract representation of the environment describing the available objects, actions, and current system state. Existing work has represented this information using natural-language descriptions, symbolic planning languages, executable code, and other structured representations. The choice of representation influences both the information available to the language model and the ease with which generated plans can be validated and executed, making environment representation an important design consideration for LLM-based task planning.

Structured representations such as JSON have become increasingly common in LLM-based planning because they encode hierarchical information through key-value pairs while remaining machine-readable and readily integrated with existing robotic software. In contrast, natural-language descriptions provide greater flexibility in expressing task context and relationships between entities. Consequently, both structured and natural-language representations have been adopted for conveying environment state to language models. However, recent work suggests that imposing restrictive formats may also influence language-model reasoning. Tam \textit{et al.}~\cite{tam2024letmespeak} showed that strict output format constraints can reduce LLM performance on reasoning tasks, highlighting a potential trade-off between structured representations and generation flexibility. Despite the increasing use of both approaches, their impact on high-level task planning performance in heterogeneous multi-robot systems remains insufficiently understood.


Language model selection represents another important design consideration in LLM-based robotic planning. In applied robotics and embedded systems, constraints on computation, memory, energy consumption, and inference latency may favour the deployment of smaller language models over larger alternatives~\cite{dawane2026onddeviceslm}. However, smaller models often exhibit weaker reasoning capabilities and may therefore require more carefully designed prompts, structured task representations, or additional validation mechanisms to achieve reliable planning performance.Xiao et al. and Mandi et al.~\cite{xiao2026probingprompt,mandi2024roco} investigated prompt design as a means of compensating for limited model capacity and improving decision quality in resource-constrained deployments. At the same time, studies have shown that imposing restrictive prompt or output formats may negatively influence language-model reasoning, suggesting a trade-off between structured machine-readable representations and generation flexibility~\cite{sahoo2024promptsurvey,tam2024letmespeak}. Collectively, these studies indicate that planning performance depends not only on the choice of language model, but also on how planning problems are represented and communicated to the model.

A variety of simulation environments have been developed to evaluate multi-robot planning and coordination algorithms. These include warehouse-specific benchmarks, industrial warehouse simulators, and general-purpose robotic simulation platforms such as NVIDIA Isaac Sim, with the choice of simulator typically depending on the planning problem under investigation~\cite{krnjaic2024warehouse,niu2025symbiosis}.

Among warehouse-specific benchmarks, the Multi-Robot Warehouse (RWARE) environment introduced by Papoudakis \textit{et al.}~\cite{papoudakis2021benchmarking} models a partially observable grid-world warehouse in which multiple robots cooperate to retrieve requested shelves and deliver them to workstations. Agents receive sparse rewards only after successfully completing a sequence of coordinated actions, making coordination difficult and providing a challenging benchmark for cooperative multi-agent reinforcement learning.

To investigate heterogeneous task allocation, Krnjaic \textit{et al.}~\cite{krnjaic2024warehouse} extended RWARE to develop the Task-Assignment Robotic Warehouse (TA-RWARE) environment. TA-RWARE introduces two heterogeneous agent roles, Automated Guided Vehicles (AGVs) and picker robots, which cooperate to complete warehouse operations. Instead of directly controlling robot motion, agents select target locations while navigation between locations is performed using a predefined path-planning heuristic. By abstracting low-level navigation, the environment enables researchers to focus on high-level task assignment and coordination in heterogeneous multi-robot systems.
